\documentclass[11pt,a4paper]{article}

% ===================== PACOTES =====================
\usepackage[utf8]{inputenc}
\usepackage[T1]{fontenc}
\usepackage[brazilian]{babel}
\usepackage[margin=2.5cm]{geometry}
\usepackage{xcolor}
\usepackage{graphicx}
\usepackage{booktabs}
\usepackage{longtable}
\usepackage{array}
\usepackage{listings}
\usepackage{hyperref}
\usepackage{fancyhdr}
\usepackage{titlesec}
\usepackage{enumitem}
\usepackage{amsmath}
\usepackage{tcolorbox}
\tcbuselibrary{skins,breakable}
\graphicspath{{images/}}


\definecolor{codebg}{HTML}{F5F5F5}
\definecolor{codeframe}{HTML}{CCCCCC}
\definecolor{linknavy}{HTML}{00599C}

\pagestyle{fancy}
\fancyhf{}
\fancyhead[L]{PNAAT Gestures ESP32}
\fancyhead[R]{\thepage}
\renewcommand{\headrulewidth}{0.4pt}


\titleformat{\section}{\Large\bfseries\color{black}}{\thesection}{1em}{}
\titleformat{\subsection}{\large\bfseries\color{black}}{\thesubsection}{1em}{}


\lstdefinestyle{cstyle}{
    backgroundcolor=\color{codebg},
    frame=single,
    rulecolor=\color{codeframe},
    basicstyle=\ttfamily\small,
    keywordstyle=\color{blue}\bfseries,
    commentstyle=\color{gray}\itshape,
    stringstyle=\color{orange},
    numbers=left,
    numberstyle=\tiny\color{gray},
    numbersep=8pt,
    breaklines=true,
    breakatwhitespace=false,
    showstringspaces=false,
    tabsize=2,
    captionpos=b,
    language=C
}
\lstset{style=cstyle}


\hypersetup{
    colorlinks=true,
    linkcolor=linknavy,
    urlcolor=linknavy,
    citecolor=linknavy,
    pdftitle={PNAAT Gestures ESP32},
    pdfauthor={Projeto PNAAT}
}


\newcommand{\projimg}[2]{%
  \IfFileExists{images/#1}{%
    \begin{figure}[h]
      \centering
      \includegraphics[width=0.8\linewidth]{#1}
      \caption{#2}
    \end{figure}
  }{%
    \begin{center}
    \fbox{\parbox{0.8\linewidth}{\centering \vspace{2cm} IMAGEM: #2 \\ \small(arquivo esperado: \texttt{docs/images/#1}) \vspace{2cm}}}
    \end{center}
  }%
}

\begin{document}


\begin{titlepage}
    \centering
    \vspace*{2cm}
    {\Huge \bfseries PNAAT Gestures ESP32}\\[0.5cm]
    {\Large Automação Residencial com Inteligência Artificial Embarcada}\\[2cm]

    \rule{\linewidth}{0.5pt}\\[0.3cm]
    {\large Sistema de reconhecimento de gestos manuais utilizando rede neural\\
    executada 100\% em um ESP32-S3, com controle de LED via PWM,\\
    relógio sincronizado por NTP, display OLED e telemetria/controle remoto via MQTT.}\\[0.3cm]
    \rule{\linewidth}{0.5pt}\\[2cm]

    \begin{tcolorbox}[colback=yellow!10,colframe=orange!70!black,width=0.85\linewidth,breakable]
        \textbf{Aviso:} Este é um projeto acadêmico de capstone/trabalho final de curso (protótipo acadêmico).
        Pode conter inconsistências, casos de borda não tratados e limitações típicas de um projeto de aprendizado —
        não é um produto pronto para produção.
    \end{tcolorbox}

    \vfill
    {\large Plataforma: ESP32-S3 (Heltec WiFi LoRa 32 V3)\\
    Framework: ESP-IDF v5.4.4\\
    Licença: MIT}\\[1cm]

    {\large \today}
\end{titlepage}

\tableofcontents
\newpage


\section{Funcionalidades}

\begin{itemize}[leftmargin=*]
    \item Reconhecimento de \textbf{5 gestos} por uma rede neural embarcada: \texttt{direita}, \texttt{esquerda}, \texttt{cima}, \texttt{baixo}, \texttt{parado}.
    \item \texttt{direita} liga o LED; \texttt{esquerda} desliga.
    \item \texttt{cima}/\texttt{baixo} aumentam/diminuem o brilho em passos de \textbf{20\%} via PWM (LEDC, 5 kHz, 10 bits = 1024 níveis).
    \item Display OLED SSD1306 128x64 exibindo, nesta ordem: horário, gesto, LED e brilho.
    \item Relógio sincronizado via \textbf{SNTP}, fuso horário de Brasília (UTC-3).
    \item Telemetria MQTT em JSON a cada ação.
    \item Controle remoto do LED via MQTT (\texttt{on}/\texttt{off}/\texttt{up}/\texttt{down}).
    \item Operação offline resiliente, com reconexão automática de Wi-Fi.
    \item Diagnóstico embutido via terminal (taxa de amostragem e contagem de pulsos no pino INT).
\end{itemize}


\section{Hardware e Pinagem}

\textbf{Placa:} Heltec WiFi LoRa 32 V3 (ESP32-S3)\\
\textbf{Sensores/periféricos:} IMU BNO085 (breakout GY-BNO08X), OLED SSD1306 128x64 onboard (endereço I2C \texttt{0x3C}), LED branco onboard (\texttt{GPIO35}).

\begin{longtable}{p{2.8cm}p{2.2cm}p{2.2cm}p{6cm}}
\toprule
\textbf{Componente} & \textbf{Sinal} & \textbf{Pino} & \textbf{Observações} \\
\midrule
\endhead
BNO085 & SDA & GPIO6 & \texttt{I2C\_NUM\_0}, 100 kHz \\
BNO085 & SCL & GPIO7 & \texttt{I2C\_NUM\_0}, 100 kHz \\
BNO085 & INT & GPIO5 & Obrigatório — sinaliza novo dado \\
BNO085 & RST & GPIO4 & Fallback: reset por software \\
BNO085 & ADD & GND & Define endereço \texttt{0x4A} \\
BNO085 & VCC & 3.3V fixo & Nunca alimentar pelo trilho Vext \\
OLED & SDA & GPIO17 & \texttt{I2C\_NUM\_1} dedicado \\
OLED & SCL & GPIO18 & \texttt{I2C\_NUM\_1} dedicado \\
OLED & RST & GPIO21 & --- \\
OLED & Alimentação & Trilho VEXT (GPIO36) & Ativo em LOW, deve ser habilitado antes da inicialização do display \\
Console & TX/RX & GPIO43/GPIO44 & UART \\
\bottomrule
\end{longtable}

\begin{tcolorbox}[colback=blue!5,colframe=linknavy,breakable]
O projeto utiliza a \textbf{nova API I2C} do ESP-IDF (\texttt{driver/i2c\_master.h}), e não a API legada.
\end{tcolorbox}

\projimg{images/1}{Foto da montagem real com jumpers visíveis}
\projimg{images/2}{Diagrama esquemático/Fritzing da ligação BNO085 $\leftrightarrow$ Heltec}


\section{Arquitetura de Software}

Framework: \textbf{ESP-IDF v5.4.4} (C/C++).

\subsection{Componentes (\texttt{components/})}
\begin{longtable}{p{3.5cm}p{9cm}}
\toprule
\textbf{Componente} & \textbf{Função} \\
\midrule
\endhead
\texttt{i2c\_config} & Inicializa os dois barramentos I2C + o trilho Vext \\
\texttt{oled\_setup} & Painel SSD1306 + LVGL via \texttt{esp\_lvgl\_port} \\
\texttt{bno085} & Driver vendorizado sobre a biblioteca SH2 (a versão do registro tinha um bug de \texttt{\%u} que quebrava o \texttt{-Werror}) \\
\texttt{edge\_impulse} & SDK do Edge Impulse + modelo treinado (\texttt{model-parameters/}, \texttt{tflite-model/}) \\
\bottomrule
\end{longtable}

\subsection{Módulos principais (\texttt{main/})}
\begin{longtable}{p{3.5cm}p{9cm}}
\toprule
\textbf{Arquivo} & \textbf{Função} \\
\midrule
\endhead
\texttt{main.c} & Orquestra o boot \\
\texttt{imu\_config.cpp} & Núcleo: PWM, janelamento de amostras, inferência, lógica de disparo de gestos, OLED, relógio, diagnóstico \\
\texttt{wifi\_ntp.c} & Wi-Fi STA + SNTP \\
\texttt{mqtt\_link.c} & Cliente MQTT com fila FreeRTOS \\
\bottomrule
\end{longtable}

\subsection{Ordem de boot}
Wi-Fi (não bloqueante) $\rightarrow$ trilho Vext $\rightarrow$ OLED $\rightarrow$ IMU $\rightarrow$ aguarda IP (até 30s) $\rightarrow$ MQTT.

\subsection{Tarefas (Tasks)}
\texttt{imu\_service} (prioridade 5), callback do sensor, \texttt{mqtt\_pub}, \texttt{clock}, \texttt{diag}, \texttt{network}, tarefa interna do LVGL.

\subsection{Partições (\texttt{partitions.csv})}
\texttt{nvs} 24K, \texttt{phy\_init} 4K, \texttt{factory} 4 MB (binário $\approx$ 1.1 MB).


\section{Pipeline de Machine Learning}

\begin{enumerate}[leftmargin=*]
    \item \textbf{Coleta de dados:} um firmware separado (\texttt{receber\_dados}) transmite CSV de \texttt{roll,pitch,yaw} a 50 Hz via o Data Forwarder do Edge Impulse. As fórmulas de conversão de quatérnion para Euler são idênticas no coletor e no firmware final:

\begin{lstlisting}[caption={Conversão quatérnion para Euler}]
roll  = atan2f(2*(qr*qi + qj*qk), 1 - 2*(qi*qi + qj*qj)) * 57.2958f;
pitch = asinf(2*(qr*qj - qk*qi)) * 57.2958f;
yaw   = atan2f(2*(qr*qk + qi*qj), 1 - 2*(qj*qj + qk*qk)) * 57.2958f;
\end{lstlisting}

    \item \textbf{Dataset:} 5 classes — \texttt{baixo}, \texttt{cima}, \texttt{direita}, \texttt{esquerda}, \texttt{parado}.
    \item \textbf{Modelo:} janela de 150 amostras $\times$ 3 eixos = 450 valores ($\approx$3s), rede quantizada em int8, compilada com o EON Compiler.
    \item \textbf{Inferência:} buffer preenchido pelo callback do sensor; assim que a janela é preenchida $\rightarrow$ \texttt{run\_classifier()}. Configurado via macros (\texttt{EI\_CLASSIFIER\_DSP\_INPUT\_FRAME\_SIZE}, \texttt{EI\_CLASSIFIER\_LABEL\_COUNT}).
    \item \textbf{Trigger anti-repetição:} uma ação só é disparada quando o rótulo vencedor muda \textbf{e} a confiança $\geq$ 60\%. Detectar \texttt{parado} rearma o trigger.
    \begin{itemize}
        \item \texttt{direita → parado → parado}: LED permanece ligado
        \item \texttt{direita → parado → direita}: nenhuma nova ação
    \end{itemize}
\end{enumerate}

\projimg{images/8}{Capturas de tela do Edge Impulse Studio — explorador de dados, acurácia/perda, matriz de confusão}
\projimg{images/3}{Gráfico de uma amostra de gesto mostrando variação de pitch/yaw}


\section{Detalhes do Treinamento (Edge Impulse)}

Esta seção documenta a configuração exata usada no Edge Impulse Studio para gerar o modelo embarcado, garantindo a reprodutibilidade do experimento.

\subsection{Configuração do Impulse (janelamento)}
\begin{table}[h]
\centering
\begin{tabular}{ll}
\toprule
\textbf{Parâmetro} & \textbf{Valor} \\
\midrule
Eixos de entrada & x, y, z (3 eixos) \\
Window size & 3.000 ms \\
Window increase (stride) & 1.000 ms \\
Frequência (Hz) & 50 Hz \\
Zero-pad data & Habilitado \\
Train on data subset & 100\% \\
Classes de saída & 5 (baixo, cima, direita, esquerda, parado) \\
\bottomrule
\end{tabular}
\caption{Parâmetros do bloco \emph{Time series data}}
\end{table}

\subsection{Processamento do sinal (DSP)}
Bloco \emph{Raw Data}, com os 3 eixos habilitados e fator de escala (\texttt{Scale axes}) igual a 1, ou seja, sem normalização adicional além do que já é feito na conversão quatérnion $\rightarrow$ Euler.

\begin{table}[h]
\centering
\begin{tabular}{ll}
\toprule
\textbf{Métrica de desempenho on-device} & \textbf{Valor} \\
\midrule
Tempo de processamento & 1 ms \\
Uso de pico de RAM & 2 KB \\
\bottomrule
\end{tabular}
\caption{Desempenho estimado do bloco DSP no dispositivo embarcado}
\end{table}

\subsection{Arquitetura da rede neural}
Rede neural densa (fully-connected), treinada e quantizada em int8 pelo EON Compiler.

\begin{table}[h]
\centering
\begin{tabular}{ll}
\toprule
\textbf{Camada} & \textbf{Detalhe} \\
\midrule
Input layer & 450 features (janela de 150 amostras $\times$ 3 eixos) \\
Dense layer & 20 neurônios \\
Dense layer & 10 neurônios \\
Output layer & 5 classes \\
\bottomrule
\end{tabular}
\caption{Arquitetura da rede neural}
\end{table}

\begin{table}[h]
\centering
\begin{tabular}{ll}
\toprule
\textbf{Parâmetro de treinamento} & \textbf{Valor} \\
\midrule
Number of training cycles & 30 \\
Learning rate & 0.0005 \\
Use learned optimizer & Desabilitado \\
Training processor & CPU \\
Model version & Quantized (int8) \\
\bottomrule
\end{tabular}
\caption{Hiperparâmetros de treinamento}
\end{table}

\subsection{Resultados (conjunto de validação)}
\begin{table}[h]
\centering
\begin{tabular}{ll}
\toprule
\textbf{Métrica} & \textbf{Valor} \\
\midrule
Acurácia & 100,0\% \\
Loss & 0,03 \\
Area under ROC Curve & 1,00 \\
Weighted average Precision & 1,00 \\
Weighted average Recall & 1,00 \\
Weighted average F1 score & 1,00 \\
\bottomrule
\end{tabular}
\caption{Métricas gerais de desempenho do modelo}
\end{table}

\begin{table}[h]
\centering
\begin{tabular}{lccccc}
\toprule
 & \textbf{BAIXO} & \textbf{CIMA} & \textbf{DIREITA} & \textbf{ESQUERDA} & \textbf{PARADO} \\
\midrule
\textbf{BAIXO}    & 100\% & 0\%  & 0\%  & 0\%  & 0\%  \\
\textbf{CIMA}     & 0\%   & 100\%& 0\%  & 0\%  & 0\%  \\
\textbf{DIREITA}  & 0\%   & 0\%  & 100\%& 0\%  & 0\%  \\
\textbf{ESQUERDA} & 0\%   & 0\%  & 0\%  & 100\%& 0\%  \\
\textbf{PARADO}   & 0\%   & 0\%  & 0\%  & 0\%  & 100\% \\
\bottomrule
\end{tabular}
\caption{Matriz de confusão (conjunto de validação) — F1 score = 1,00 em todas as classes}
\end{table}

\begin{tcolorbox}[colback=yellow!10,colframe=orange!70!black,breakable]
\textbf{Observação sobre os resultados:} a acurácia de 100\% no conjunto de validação é um indício positivo, mas também pode
refletir um dataset relativamente pequeno ou pouco diverso (poucos usuários/condições de coleta). Recomenda-se validar o
modelo com dados coletados em sessões e condições diferentes das usadas no treinamento, para confirmar sua capacidade de
generalização antes de considerá-lo definitivo.
\end{tcolorbox}

\projimg{images/4}{Tela de treinamento do classificador com matriz de confusão}


\section{Controle de LED via PWM}

LEDC: canal 0, timer 0, modo baixa velocidade, 5 kHz, 10 bits (duty 0-1023).

\begin{table}[h]
\centering
\begin{tabular}{ll}
\toprule
\textbf{Gesto} & \textbf{Ação} \\
\midrule
\texttt{direita} & LED LIGADO \\
\texttt{esquerda} & LED DESLIGADO \\
\texttt{cima} & +20\% de brilho \\
\texttt{baixo} & -20\% de brilho \\
\bottomrule
\end{tabular}
\end{table}

\begin{lstlisting}[caption={Cálculo do duty cycle}]
duty = (brightness_percent * 1023) / 100;  // OFF = duty 0
\end{lstlisting}

\projimg{led-brilho.jpg}{Sequência de fotos do LED branco a 0\%, 40\% e 80\% de brilho}


\section{Display OLED e Relógio NTP}

4 linhas, exibidas nesta ordem:
\begin{enumerate}[leftmargin=*]
    \item \textbf{Horário} — relógio no formato \texttt{HH:MM:SS}
    \item \textbf{Gesto} — último gesto detectado
    \item \textbf{LED} — estado ligado/desligado
    \item \textbf{Brilho} — porcentagem atual
\end{enumerate}

Atualiza a cada janela classificada ou comando remoto (mostra a origem, ex.: \texttt{remoto:on}). Antes da sincronização, o horário exibe \texttt{--:--:--}.

Sem RTC com backup de bateria, a placa sempre inicializa na época Unix. Uma vez conectada ao Wi-Fi, o SNTP sincroniza com \texttt{pool.ntp.org} e \texttt{a.st1.ntp.br}, fuso horário \texttt{<-03>3} (Brasília, sem horário de verão).


\section{Conectividade IoT / MQTT}

\textbf{Broker:} \texttt{broker.hivemq.com:1883}

\subsection{Publicação}
Tópico \texttt{testtopic/pnaat/status} (QoS 1), ao conectar, a cada gesto e a cada comando remoto:

\begin{lstlisting}[language=json,caption={Exemplo de payload de telemetria}]
{
  "source": "cima",
  "led": true,
  "brightness": 80,
  "time": "21:03:12"
}
\end{lstlisting}

\subsection{Assinatura}
Tópico \texttt{testtopic/pnaat/cmd}, payloads: \texttt{on}, \texttt{off}, \texttt{up}, \texttt{down}.

\subsection{Semântica híbrida}
Última escrita vence — um comando remoto é aplicado imediatamente, mas os gestos locais continuam funcionando depois.

\subsection{Resiliência}
A publicação passa por uma fila FreeRTOS fora do caminho crítico do sensor; sem Wi-Fi/internet tudo continua funcionando localmente, com reconexão automática. Credenciais configuráveis via \texttt{menuconfig} (seção "Configuracao Wi-Fi / MQTT (PNAAT)").

\projimg{mqtt-telemetria.png}{Captura de tela da telemetria JSON chegando em um cliente MQTT}
\projimg{mqtt-comando-on.jpg}{Foto do LED branco acendendo enquanto o console envia o comando \texttt{on}}


\section{Diagnóstico Embutido}

A tarefa \texttt{diag} imprime a cada 5s:

\begin{lstlisting}[language={}]
[diag] X samples/s | Y INT pulses/s | windows: Z
\end{lstlisting}

Esperado: $\approx$50 amostras/s e $\approx$100 pulsos/s.

\begin{table}[h]
\centering
\begin{tabular}{lll}
\toprule
\textbf{Amostras} & \textbf{Pulsos INT} & \textbf{Diagnóstico} \\
\midrule
0 & 0 & Problema de fiação no INT \\
0 & OK & Problema de software \\
OK & OK & Sistema saudável \\
\bottomrule
\end{tabular}
\end{table}


\section{Como Compilar e Gravar}

\begin{lstlisting}[language=bash,caption={Compilação e gravação}]
source ~/esp/v5.4.4/esp-idf/export.sh
idf.py menuconfig      # preencha SSID/senha na secao PNAAT
idf.py -p ttyUSB0 build flash monitor
\end{lstlisting}

Console MQTT opcional em \texttt{tools/}:

\begin{lstlisting}[language=bash]
./tools/pnaat_mqtt.sh
\end{lstlisting}

Terminal interativo mostrando telemetria formatada (\texttt{time | source | LED | brightness}) e enviando comandos.

\textbf{Configurações relevantes no \texttt{sdkconfig}:} otimização \texttt{-Os}, \texttt{CONFIG\_PARTITION\_TABLE\_CUSTOM}, flash de 8 MB, \texttt{LV\_FONT\_MONTSERRAT\_12} habilitado.


\section{Desafios Técnicos Resolvidos}

\begin{enumerate}[leftmargin=*]
    \item Erro de link: \texttt{imu\_config\_init()} declarada/chamada mas nunca definida.
    \item Migração da API I2C legada para \texttt{driver/i2c\_master.h}.
    \item Bug no componente do registro \texttt{rinku404/bno085} (\texttt{\%u} quebrava o \texttt{-Werror}) — corrigido vendorizando o componente.
    \item Conflito entre um callback de flush customizado e o callback interno do \texttt{esp\_lvgl\_port}.
    \item Estouro da partição padrão de 1 MB — corrigido com tabela de partições customizada de 4 MB e correção do tamanho de flash no \texttt{sdkconfig} (2 MB configurado vs. 8 MB real).
    \item Binário inchado por causa do \texttt{-Og} — trocado para \texttt{-Os}, suprimindo falsos positivos de \texttt{-Wuninitialized} vindos dos headers do SDK do Edge Impulse.
    \item Condição de corrida no MQTT (erro DNS \texttt{getaddrinfo 202}) — corrigido iniciando o MQTT somente após o \texttt{GOT\_IP}.
    \item Falhas intermitentes na IMU — causa raiz: um jumper solto no pino INT, encontrado via diagnóstico.
    \item Depreciações do GCC 14 (incremento de \texttt{volatile}, \texttt{std::is\_pod}).
\end{enumerate}


\section{Estrutura do Repositório}

\begin{lstlisting}[language={}, basicstyle=\ttfamily\footnotesize]
Projeto_finalPNAAT/
|-- build/
|-- components/
|   |-- i2c_config/
|   |-- oled_setup/
|   `-- bno085/              # vendorizado
|-- docs/
|   |-- images/               # imagens deste documento
|   `-- documentacao.tex      # este arquivo
|-- main/
|   |-- main.c
|   |-- imu_config.cpp
|   |-- imu_config.h
|   |-- wifi_ntp.c
|   |-- wifi_ntp.h
|   |-- mqtt_link.c
|   |-- mqtt_link.h
|   |-- Kconfig.projbuild
|   `-- idf_component.yml
|-- managed_components/
|-- tools/
|   |-- pnaat_mqtt.sh
|   `-- pnaat_mqtt.py
|-- CMakeLists.txt
|-- LICENSE
|-- README.md
|-- dependencies.lock
|-- partitions.csv
|-- sdkconfig
`-- sdkconfig.old
\end{lstlisting}


\section{Tecnologias Utilizadas}

\begin{itemize}[leftmargin=*]
    \item ESP-IDF v5.4.4
    \item FreeRTOS
    \item Edge Impulse + TensorFlow Lite Micro (EON)
    \item LVGL 9.2
    \item esp\_lvgl\_port
    \item MQTT (esp-mqtt)
    \item SNTP/LwIP
    \item Driver LEDC
    \item Nova API \texttt{i2c\_master}
\end{itemize}


\section{Créditos e Agradecimentos}

Curso PNAAT $\cdot$ Plataforma Edge Impulse $\cdot$ Heltec $\cdot$ Comunidade Espressif


\section{Licença}

Este projeto está licenciado sob a Licença MIT — consulte o arquivo \texttt{LICENSE} para mais detalhes.

\vspace{1cm}
\begin{center}
\textit{\textbf{100\% IA de borda (Edge AI)} — a classificação acontece no microcontrolador, sem dependência de nuvem.}
\end{center}

\end{document}
